\chapter{Cơ sở lý thuyết}
\label{ch:background}

Chương này trình bày các khái niệm nền tảng được sử dụng trong toàn bộ báo cáo: định nghĩa
hình thức của phép scan (\cref{sec:scan-def}), mô hình phân tích thuật toán song song
work--span (\cref{sec:workspan}), các định luật giới hạn tăng tốc và mô hình roofline
(\cref{sec:laws}), chuẩn lập trình OpenMP (\cref{sec:openmp}) và các đặc điểm kiến trúc CPU đa
nhân ảnh hưởng trực tiếp tới hiệu năng (\cref{sec:arch}).

\section{Phép tính tổng tiền tố}
\label{sec:scan-def}

\begin{definition}[Inclusive scan]
Cho phép toán hai ngôi \(\oplus\) có tính kết hợp trên tập \(X\) và dãy
\(x = [x_0, x_1, \ldots, x_{n-1}] \in X^n\). \emph{Inclusive scan} của \(x\) là dãy
\[
  \operatorname{scan}(x) = [\,x_0,\; x_0 \oplus x_1,\; \ldots,\; x_0 \oplus x_1 \oplus \cdots \oplus x_{n-1}\,].
\]
\end{definition}

\begin{definition}[Exclusive scan / prescan]
Nếu \(\oplus\) có phần tử đơn vị \(I\), \emph{exclusive scan} (Blelloch gọi là
\emph{prescan}~\cite{blelloch1990}) của \(x\) là
\[
  \operatorname{prescan}(x) = [\,I,\; x_0,\; x_0 \oplus x_1,\; \ldots,\; x_0 \oplus \cdots \oplus x_{n-2}\,].
\]
\end{definition}

\begin{example}
Với \(\oplus = +\), \(I = 0\) và \(x = [3, 1, 7, 0, 4, 1, 6, 3]\):
\[
  \operatorname{scan}(x) = [3, 4, 11, 11, 15, 16, 22, 25], \qquad
  \operatorname{prescan}(x) = [0, 3, 4, 11, 11, 15, 16, 22].
\]
\end{example}

Hai dạng có thể chuyển đổi cho nhau với chi phí \(\Oh(n)\) và hoàn toàn song song:
\(\operatorname{scan}(x)_i = \operatorname{prescan}(x)_i \oplus x_i\), và ngược lại
\(\operatorname{prescan}(x)_i = \operatorname{scan}(x)_{i-1}\) với \(\operatorname{prescan}(x)_0 = I\).
Báo cáo này dùng inclusive scan làm đầu ra chuẩn cho mọi phiên bản; thuật toán Blelloch tự
nhiên sinh ra exclusive scan và được chuyển đổi ở bước cuối.

\paragraph{Vai trò của tính kết hợp.} Mọi thuật toán scan song song đều dựa trên việc
\emph{nhóm lại} các phép toán theo thứ tự khác với thứ tự tuần tự, ví dụ
\((x_0 \oplus x_1) \oplus (x_2 \oplus x_3)\) thay vì \(((x_0 \oplus x_1) \oplus x_2) \oplus x_3\).
Điều này chỉ hợp lệ khi \(\oplus\) kết hợp. Tính giao hoán \emph{không} cần thiết: các thuật toán
giữ nguyên thứ tự trái--phải của toán hạng. Do đó scan áp dụng được cho nhiều phép toán:
cộng, nhân, \(\max\), \(\min\), AND/OR logic, nhân ma trận, hợp các hàm affine
\(f(y) = ay + b\) (dùng để giải hệ truy hồi tuyến tính bậc nhất)~\cite{blelloch1990}.

\paragraph{Vấn đề với số thực.} Phép cộng dấu phẩy động IEEE~754 chỉ \emph{gần} kết hợp:
\((a + b) + c\) và \(a + (b + c)\) có thể khác nhau ở bit cuối do làm tròn. Thuật toán song song
vì vậy cho kết quả hơi khác bản tuần tự; sai số tích lũy của cách nhóm theo cây thậm chí thường
\emph{nhỏ hơn} cách cộng tuần tự. Báo cáo dùng số nguyên 64 bit để có thể so sánh chính xác
từng bit.

\paragraph{Thuật toán tuần tự.} Cài đặt tuần tự rất đơn giản (\cref{alg:seq}): đúng \(n-1\)
phép cộng, đây là số phép toán tối thiểu cần thiết, đồng thời là mốc để đánh giá các thuật toán
song song. Tuy nhiên vòng lặp có \emph{phụ thuộc mang qua vòng lặp} (loop-carried dependency):
bước \(i\) cần kết quả bước \(i-1\), nên không thể đơn giản thêm \omp{parallel for}.

\begin{algorithm}[H]
\caption{Inclusive scan tuần tự}
\label{alg:seq}
\KwIn{mảng \(x[0..n-1]\)}
\KwOut{mảng \(y[0..n-1]\), \(y_i = \sum_{j \le i} x_j\)}
\(s \gets 0\)\;
\For{\(i \gets 0\) \KwTo \(n-1\)}{
  \(s \gets s + x[i]\)\;
  \(y[i] \gets s\)\;
}
\end{algorithm}

\section{Mô hình phân tích thuật toán song song}
\label{sec:workspan}

\subsection{Mô hình PRAM}

Mô hình \emph{Parallel Random Access Machine} (PRAM) gồm \(p\) bộ xử lý đồng bộ cùng truy cập
một bộ nhớ chung; mỗi bước, mọi bộ xử lý thực hiện một lệnh. Biến thể EREW (Exclusive Read,
Exclusive Write) cấm hai bộ xử lý cùng đọc hoặc cùng ghi một ô nhớ trong một bước. Blelloch
phân tích các thuật toán scan trên EREW PRAM~\cite{blelloch1990}. PRAM bỏ qua chi phí đồng bộ,
độ trễ và băng thông bộ nhớ -- những yếu tố lại chiếm ưu thế trên phần cứng thật, như ta sẽ
thấy ở \cref{ch:experiments}.

\subsection{Work, span và định lý Brent}

Một thuật toán song song được mô tả bằng một \emph{đồ thị phụ thuộc} có hướng không chu trình
(DAG), mỗi nút là một phép toán, mỗi cạnh là một phụ thuộc dữ liệu. Hai đại lượng đặc trưng:
\begin{itemize}
  \item \textbf{Work} \(\Work(n)\): tổng số nút -- thời gian chạy trên một bộ xử lý, \(T_1\).
  \item \textbf{Span} (hay \emph{depth}) \(\Span(n)\): độ dài đường đi dài nhất trong DAG --
        thời gian chạy khi có vô hạn bộ xử lý, \(T_\infty\).
\end{itemize}
Tỉ số \(\Work/\Span\) gọi là \emph{độ song song} (parallelism) -- số bộ xử lý tối đa có thể tận
dụng hiệu quả. Với \(p\) bộ xử lý, hiển nhiên \(T_p \ge \max(\Work/p,\; \Span)\).

\begin{theorem}[Brent~\cite{brent1974}]
Một thuật toán có work \(\Work\) và span \(\Span\) có thể được thực thi trên \(p\) bộ xử lý trong
thời gian
\begin{equation}
  T_p \;\le\; \frac{\Work}{p} + \Span.
  \label{eq:brent}
\end{equation}
\end{theorem}

Định lý Brent cho thấy hai mục tiêu thiết kế: giảm \(\Span\) để có nhiều song song, và
\emph{không tăng} \(\Work\) so với thuật toán tuần tự tốt nhất.

\begin{definition}[Work-efficient]
Thuật toán song song gọi là \emph{work-efficient} nếu \(\Work(n) = \Oh(T_{\text{seq}}(n))\),
trong đó \(T_{\text{seq}}\) là thời gian của thuật toán tuần tự tốt nhất.
\end{definition}

Work-efficiency đặc biệt quan trọng khi \(p \ll n\) -- trường hợp luôn gặp trên CPU (\(p\) cỡ
10, \(n\) cỡ \(10^8\)). Khi đó số hạng \(\Work/p\) trong~\eqref{eq:brent} chiếm ưu thế
tuyệt đối; một thuật toán có \(\Work = \Oh(n \log n)\) làm nhiều việc hơn bản tuần tự
\(\log_2 n\) lần (khoảng 27 lần với \(n = 2^{27}\)), nên cần ít nhất \(\log_2 n\) bộ xử lý chỉ để
hòa vốn. Đây là lý do thuật toán Hillis--Steele, dù có span tốt nhất, lại chậm nhất trong thực
nghiệm (\cref{sec:exp-strong}).

\subsection{Speedup và hiệu suất}

Với \(T_{\text{seq}}\) là thời gian của \emph{bản tuần tự tốt nhất} (không phải bản song song
chạy 1 luồng) và \(T_p\) là thời gian với \(p\) luồng:
\begin{equation}
  S_p = \frac{T_{\text{seq}}}{T_p}, \qquad E_p = \frac{S_p}{p}.
  \label{eq:speedup}
\end{equation}
\(S_p\) là \emph{speedup} (hệ số tăng tốc), \(E_p\) là \emph{hiệu suất song song} (efficiency).
Lý tưởng \(S_p = p\), \(E_p = 1\). Báo cáo luôn so sánh với bản tuần tự để tránh ``thổi phồng''
speedup: một thuật toán song song chậm khi chạy 1 luồng có thể có \(T_1/T_p\) đẹp nhưng vẫn chậm
hơn bản tuần tự.

\section{Giới hạn của tăng tốc}
\label{sec:laws}

\subsection{Định luật Amdahl và Gustafson}

Nếu một phần \(f\) của chương trình buộc phải chạy tuần tự, định luật
Amdahl~\cite{amdahl1967} cho
\begin{equation}
  S_p \;\le\; \frac{1}{f + (1-f)/p} \;\xrightarrow{p\to\infty}\; \frac{1}{f}.
  \label{eq:amdahl}
\end{equation}
Với scan chia khối (\cref{sec:alg-blocked}), phần tuần tự là việc quét \(p\) tổng khối, có chi
phí \(\Oh(p)\) so với \(\Oh(n)\) của toàn bộ, tức \(f \approx p/n\) -- không đáng kể với \(n\) lớn.
Như vậy theo Amdahl, scan gần như song song hóa hoàn toàn được. Định luật
Gustafson~\cite{gustafson1988} xét bài toán có kích thước tăng theo \(p\) (weak scaling) và cho
\(S_p = p - f(p-1)\), lạc quan hơn nữa.

Tuy nhiên, cả hai định luật chỉ xét \emph{số phép tính}. Như phần tiếp theo chỉ ra, với scan
yếu tố quyết định lại là \emph{lượng dữ liệu phải di chuyển}.

\subsection{Mô hình roofline}
\label{sec:roofline}

Mô hình roofline~\cite{williams2009} đặc trưng một chương trình bằng \emph{cường độ tính toán}
(arithmetic intensity) \(I\) = số phép toán / số byte trao đổi với bộ nhớ chính, và chặn trên
hiệu năng đạt được:
\begin{equation}
  P \;\le\; \min\bigl(P_{\text{peak}},\; I \times \beta\bigr),
  \label{eq:roofline}
\end{equation}
trong đó \(P_{\text{peak}}\) là hiệu năng tính toán đỉnh và \(\beta\) là băng thông bộ nhớ.
Chương trình có \(I\) nhỏ nằm ở vùng ``mái nghiêng'' -- \emph{giới hạn bởi băng thông}
(memory-bound): thêm nhân cũng không nhanh hơn khi băng thông đã bão hòa.

Với scan trên số nguyên 64 bit, mỗi phần tử cần ít nhất 1 phép cộng, đọc 8 byte và ghi 8 byte,
nên
\[
  I_{\text{scan}} \le \frac{1}{16}\ \text{phép cộng/byte} \approx 0{,}06.
\]
Một nhân CPU hiện đại thực hiện được hàng tỉ phép cộng mỗi giây, trong khi băng thông DRAM của
máy tính xách tay chỉ khoảng 15--60~GB/s. Tại \(I = 0{,}06\), băng thông \(15\)~GB/s chỉ ``nuôi''
được khoảng \(10^9\) phần tử/s -- mức mà \emph{một} nhân đã đạt tới. Hệ quả quan trọng:
\begin{quote}
\itshape Với dữ liệu lớn hơn bộ nhớ đệm, thời gian scan bị chặn dưới bởi
\(T \ge 16n/\beta\); speedup tối đa đạt được xấp xỉ \(\beta_{\text{đa luồng}}/\beta_{\text{đơn
luồng}}\), chứ không phải \(p\).
\end{quote}
Để có mốc băng thông thực tế, báo cáo đo kernel \emph{copy} song song (\(y_i = x_i\), cũng
16~byte/phần tử, tương tự STREAM-copy~\cite{mccalpin1995}): đây là ``trần'' mà không thuật toán
scan nào có thể vượt qua khi dữ liệu nằm ngoài cache.

\section{Chuẩn lập trình OpenMP}
\label{sec:openmp}

OpenMP~\cite{openmp52,chapman2008} là chuẩn lập trình song song cho bộ nhớ chia sẻ, gồm các chỉ
thị biên dịch (\code{\#pragma omp}), hàm thư viện (\code{omp\_*}) và biến môi trường
(\code{OMP\_*}). Ưu điểm lớn của OpenMP là cho phép song song hóa \emph{tăng dần} một chương
trình tuần tự: bỏ các chỉ thị đi, chương trình vẫn đúng.

\subsection{Mô hình fork--join}

Chương trình bắt đầu với một luồng chính. Gặp vùng \omp{parallel}, luồng chính tạo (\emph{fork})
một nhóm \(p\) luồng cùng thực thi khối lệnh; cuối vùng, các luồng đồng bộ và hợp lại
(\emph{join}) thành một (\cref{fig:forkjoin}). Thư viện \code{libgomp} giữ lại nhóm luồng giữa
các vùng song song nên chi phí fork--join chỉ vài micro giây, nhưng vẫn đáng kể khi lặp lại nhiều
lần với dữ liệu nhỏ.

\begin{figure}[H]
\centering
\begin{tikzpicture}[x=1cm,y=0.55cm,>=Stealth,thick]
  \draw[->] (0,0) -- (2,0) node[midway,above]{\small tuần tự};
  \foreach \y in {1.5,0.5,-0.5,-1.5}{
    \draw[husblue] (2,0) -- (2.6,\y) -- (6.4,\y) -- (7,0);
  }
  \node[husblue] at (4.5,2.3) {\small vùng \code{parallel} (\(p = 4\) luồng)};
  \node at (2,-2.6) {\small fork};
  \node at (7,-2.6) {\small join (barrier)};
  \draw[->] (7,0) -- (9,0) node[midway,above]{\small tuần tự};
\end{tikzpicture}
\caption{Mô hình fork--join của OpenMP.}
\label{fig:forkjoin}
\end{figure}

\subsection{Các chỉ thị sử dụng trong đề tài}

\begin{table}[H]
\centering
\caption{Các cấu trúc OpenMP được dùng để cài đặt scan.}
\label{tab:omp-constructs}
\renewcommand{\arraystretch}{1.2}
\begin{tabularx}{\textwidth}{>{\ttfamily\small}p{4.6cm}X}
\toprule
\textnormal{\textbf{Cấu trúc}} & \textbf{Ý nghĩa} \\
\midrule
parallel & Tạo nhóm luồng; mọi luồng thực thi khối lệnh. \\
for schedule(static) & Chia vòng lặp thành \(p\) đoạn liên tiếp bằng nhau; rào chắn ngầm định
ở cuối. \\
for \dots{} nowait & Bỏ rào chắn ngầm định ở cuối vòng \code{for}. \\
barrier & Rào chắn tường minh: mọi luồng chờ nhau tại điểm này. \\
single & Chỉ một luồng thực thi khối lệnh; các luồng khác chờ ở cuối. \\
simd & Yêu cầu vector hóa vòng lặp (AVX2: 4 số int64 mỗi lệnh). \\
reduction(+:s) & Mỗi luồng có bản sao riêng của \code{s}, cộng dồn ở cuối. \\
reduction(inscan,+:s) + scan inclusive(s) & Chỉ thị scan chuẩn của OpenMP 5.0: trình biên dịch
tự sinh thuật toán scan song song. \\
omp\_get\_thread\_num(), omp\_get\_num\_threads() & Chỉ số luồng và số luồng trong nhóm -- dùng
để tự chia dữ liệu theo khối. \\
omp\_get\_wtime() & Đồng hồ thời gian thực độ phân giải cao, dùng để đo hiệu năng. \\
\bottomrule
\end{tabularx}
\end{table}

\subsection{Chỉ thị \texttt{scan} của OpenMP 5.0}

Từ phiên bản 5.0 (2018), OpenMP hỗ trợ trực tiếp phép scan qua mệnh đề
\code{reduction(inscan, op: var)} trên vòng lặp và chỉ thị \code{\#pragma omp scan
inclusive(var)} (hoặc \code{exclusive}), chia thân vòng lặp thành pha ``cập nhật'' trước chỉ
thị và pha ``sử dụng'' sau chỉ thị (\cref{lst:omp-scan-intro}). Lập trình viên chỉ khai báo ý
định; việc chọn thuật toán thuộc về trình biên dịch.

\begin{lstlisting}[caption={Inclusive scan bằng chỉ thị chuẩn OpenMP 5.0.},label={lst:omp-scan-intro}]
elem_t s = 0;
#pragma omp parallel for reduction(inscan, +:s)
for (size_t i = 0; i < n; i++) {
    s += in[i];                 // pha cập nhật
    #pragma omp scan inclusive(s)
    out[i] = s;                 // pha sử dụng
}
\end{lstlisting}

Chuẩn quy định một số ràng buộc: không được dùng \code{schedule} cùng \code{reduction(inscan)},
và thân vòng lặp có đúng một chỉ thị \code{scan}. Hiệu năng của cách tiếp cận này được đánh giá
ở \cref{ch:experiments}.

\section{Kiến trúc CPU đa nhân và các yếu tố ảnh hưởng hiệu năng}
\label{sec:arch}

\subsection{Phân cấp bộ nhớ}

CPU hiện đại có nhiều cấp bộ nhớ đệm giữa nhân và DRAM: L1 (vài chục KiB, độ trễ khoảng 1~ns),
L2 (vài trăm KiB đến vài MiB, khoảng 3--5~ns) và L3 dùng chung (vài MiB đến vài chục MiB,
khoảng 10--20~ns), trong khi DRAM có độ trễ khoảng 80--100~ns~\cite{hennessy2019}. Dữ liệu
được trao đổi theo \emph{dòng cache} (cache line) 64~byte. Hai hệ quả với scan:
\begin{itemize}
  \item \textbf{Truy cập tuần tự là tốt nhất}: bộ nạp trước (prefetcher) phần cứng nhận ra mẫu
        truy cập liên tiếp và nạp trước dữ liệu, che giấu độ trễ. Truy cập với bước nhảy lớn
        (như thuật toán Blelloch ở các tầng cao của cây) chỉ dùng 8 byte trong mỗi dòng 64 byte
        được nạp, lãng phí 7/8 băng thông.
  \item \textbf{Dữ liệu vừa cache thì nhanh hơn nhiều}: khi tổng kích thước mảng vào và ra nhỏ
        hơn L3, băng thông khả dụng cao hơn DRAM nhiều lần và mỗi nhân có L1/L2 riêng, nên
        speedup có thể gần tuyến tính.
\end{itemize}

\subsection{False sharing}

Khi hai luồng ghi vào hai biến \emph{khác nhau} nằm trên cùng một dòng cache, giao thức nhất
quán cache (MESI) buộc dòng đó liên tục bị chuyển qua lại giữa các nhân -- hiện tượng
\emph{false sharing}. Trong scan chia khối, mảng \code{sums[p]} chứa tổng khối của các luồng là
ứng viên điển hình: 8 phần tử int64 liên tiếp nằm cùng một dòng 64 byte. Cách khắc phục là
\emph{đệm} (padding) mỗi phần tử ra đủ 64 byte -- vấn đề tương tự \emph{bank conflict} trong
shared memory của GPU, mà Harris và cộng sự giải quyết bằng macro
\code{CONFLICT\_FREE\_OFFSET}~\cite{harris2007}.

\subsection{CPU lai: P-core và E-core}

Từ thế hệ thứ 12, Intel dùng kiến trúc lai: \emph{P-core} hiệu năng cao (có Hyper-Threading,
2 luồng/nhân) và \emph{E-core} tiết kiệm điện (1 luồng/nhân, xung nhịp và cache thấp hơn).
CPU thực nghiệm Intel Core i5-1335U có 2 P-core và 8 E-core, tổng 12 luồng
(\cref{sec:exp-env}). Điều này ảnh hưởng tới song song hóa:
\begin{itemize}
  \item \textbf{Mất cân bằng tải}: chia đều dữ liệu (\code{schedule(static)}) khiến luồng chạy
        trên E-core xong chậm hơn; tại mỗi barrier mọi luồng phải chờ luồng chậm nhất.
  \item \textbf{Hyper-Threading}: hai luồng trên cùng P-core chia sẻ đơn vị thực thi và cache
        L1/L2; với bài toán memory-bound, luồng thứ hai gần như không thêm băng thông.
  \item \textbf{Số luồng tối ưu} vì vậy thường nhỏ hơn số luồng logic.
\end{itemize}

\subsection{First-touch và quản lý trang}

Hệ điều hành cấp phát trang nhớ vật lý tại lần ghi \emph{đầu tiên} (first-touch). Nếu đo ngay
lần đầu dùng một mảng mới, kết quả sẽ gồm cả chi phí lỗi trang (page fault). Vì vậy mọi mảng
trong thực nghiệm đều được khởi tạo song song trước khi đo, kèm vài lần chạy làm nóng (warm-up).
Trên hệ thống NUMA nhiều socket, first-touch còn quyết định dữ liệu nằm ở nút nhớ nào; máy thực
nghiệm chỉ có một socket nên yếu tố này không đáng kể.

\section{Ứng dụng của scan}
\label{sec:applications}

\Cref{tab:apps} tóm tắt các ứng dụng tiêu biểu được Blelloch~\cite{blelloch1990} và Harris và
cộng sự~\cite{harris2007} trình bày.

\begin{table}[H]
\centering
\caption{Một số ứng dụng tiêu biểu của scan.}
\label{tab:apps}
\renewcommand{\arraystretch}{1.2}
\begin{tabularx}{\textwidth}{>{\raggedright\arraybackslash}p{3.5cm}>{\centering\arraybackslash}p{1.9cm}X}
\toprule
\textbf{Ứng dụng} & \textbf{Phép \(\oplus\)} & \textbf{Ý tưởng} \\
\midrule
Stream compaction & \(+\) & Scan trên mảng cờ 0/1 cho ra vị trí đích của phần tử được giữ lại
(\cref{sec:exp-app}). \\
Radix sort & \(+\) & Mỗi lượt theo một bit: scan trên cờ bit để tính vị trí mới (phép
\emph{split}). \\
Quicksort song song & \(+\) & Phân hoạch quanh chốt bằng scan; dùng \emph{segmented scan} cho
nhiều đoạn cùng lúc. \\
Line-of-sight & \(\max\) & Một điểm nhìn thấy được nếu góc của nó lớn hơn max-prescan của các góc
phía trước. \\
Hệ truy hồi tuyến tính & hợp hàm affine & \(x_i = a_i x_{i-1} + b_i\) giải bằng scan với phép
hợp \((a,b)\circ(c,d) = (ac,\, bc + d)\). \\
Cộng số lớn & lan truyền nhớ & Lan truyền bit nhớ giữa các chữ số. \\
Summed-area table & \(+\) & Scan theo hàng rồi theo cột; lọc ảnh với cửa sổ bất kỳ trong
\(\Oh(1)\). \\
Cấp phát động & \(+\) & Mỗi luồng yêu cầu \(k_i\) phần tử; prescan cho ra địa chỉ bắt đầu. \\
\bottomrule
\end{tabularx}
\end{table}
